\documentclass[10pt,a4paper,sans]{moderncv}

% moderncv themes
\moderncvstyle{classic}
\moderncvcolor{burgundy}
\renewcommand*{\namefont}{\fontsize{20}{20}\mdseries\upshape}
\renewcommand*{\addressfont}{\fontsize{8}{8}\mdseries\upshape}
\nopagenumbers{}

% character encoding
\usepackage[utf8]{inputenc}

% adjust the page margins
\usepackage[scale=0.82]{geometry}
\setlength{\hintscolumnwidth}{3.3cm}

\usepackage{enumitem}

% personal data
\name{Nicolas}{Assouad}
\homepage{nicolasassouad.com}
\social[linkedin]{nicolas-assouad-02548412a/}
\social[github]{fondation451}


\begin{document}


\makecvtitle


\section{Skills}
\cvitem{Web3}{Blockchain, EVM, Solidity, Starknet, Stellar}
\cvitem{Language}{TS, JS, Rust, OCaml, Python, C/C++}
\cvitem{Framework}{React, Next.js, Vue, Svelte, Angular}
\cvitem{Database}{MongoDb, PostgreSQL}
\cvitem{Tool}{Git, Docker}
\cvitem{DevOps}{GCP, AWS, Vercel, Deno Deploy, Ansible}
\cvitem{Soft Skill}{Management, Mentoring, Autonomy, Initiative, Communication}

\section{Professional experiences}

\cventry{09/2024 -- Present}{\href{https://www.spiko.io/}{Senior Software Engineer} (TS, Effect, React, Blockchain, EVM, Starknet, Stellar, Git)}{\href{https://www.spiko.io/}{Spiko}}{Paris}{}{
  \begin{itemize}[label=\textbullet]
    \item Working on the all aspects of the platform
  \end{itemize}
}
\cventry{10/2023 -- 09/2024}{\href{https://www.lucca-hr.com/}{Senior Software Engineer} (C\#, TS, Angular, SQL, Git)}{\href{https://www.lucca-hr.com/}{Lucca}}{Paris}{}{
  \begin{itemize}[label=\textbullet]
    \item Transfer of Compwise features inside Lucca's products
  \end{itemize}
}
\cventry{10/2022 -- 10/2023}{\href{https://www.rhmatin.com/paie/remuneration/sirh-lucca-acquiert-compwise-pour-disposer-d-une-touche-comp-ben.html}{Co-founder \& CTO} (TS, NodeJs, GCP run, MongoDb, React, Git, Docker, Entrepreneurship, User interview, UX research)}{\href{https://www.rhmatin.com/paie/remuneration/sirh-lucca-acquiert-compwise-pour-disposer-d-une-touche-comp-ben.html}{Compwise (bought by Lucca)}}{Paris}{}{
  Compwise is a platform for managing compensation policy.
  \begin{itemize}[label=\textbullet]
    \item Fast execution, ideation to MVP: 1 month.
    \item 100+ user interviews done with HRs, CFOs, and CEOs.
    \item Design, architecture and implementation of the different versions of the platform.
    \item 100 000+ employee salary data collected to build our benchmark.
    \item 200 companies in free trial / 5 companies in paying plans.
    \item Bought by Lucca (\emph{https://www.rhmatin.com/paie/remuneration/sirh-lucca-acquiert-compwise-pour-disposer-d-une-touche-comp-ben.html})
  \end{itemize}
}
\cventry{09/2022 -- 12/2022}{\href{https://www.joinef.com/}{Founder in Residence} (Entrepreneurship, User interview, UX research)}{\href{https://www.joinef.com/}{Entrepreneur First}}{Paris}{}{
  \begin{itemize}[label=\textbullet]
    \item €100k investment to support the growth of Compwise
  \end{itemize}
}
\cventry{05/2022 -- 08/2022}{\href{https://www.joko.com/en}{Senior Software Engineer} (TS, NodeJs, AWS, Lambda, DynamoDb, React, Git)}{\href{https://www.joko.com/en}{Joko}}{Paris}{}{
  \begin{itemize}[label=\textbullet]
    \item Internationalization of the chrome extension and the web app.
    \item Conducting fit and technical interviews for new engineers.
  \end{itemize}
}
\cventry{07/2021 -- 05/2022}{\href{https://www.mindmesh.com/}{Senior Software Engineer} (TS, NodeJs, Docker, Bash, Git, PostgreSQL, React, Python)}{\href{https://www.mindmesh.com/}{Mindmesh (YC S21)}}{Paris}{}{
  Mindmesh is an early stage international startup based in Boston and Paris which develop a productivity webapp for product managers.
  \begin{itemize}[label=\textbullet]
    \item User research with the CEO and the CTO.
    \item Implementation of successive MVPs (3 pivots).
    \item Migration of the drive/slack/jira integrations from Python to Typescript fully typed.
    \item Migration of the rich text editor framework from Draft.js to Plate and open source contributions.
    \item Data visualization with D3.js.
    \item Enforce good practice (introduce strict Typescript option and linter in an existing codebase, code review).
  \end{itemize}
}
\cventry{09/2020 -- 06/2021}{\href{https://eig.etalab.gouv.fr/defis/label/}{Software Engineer} (TS, NodeJs, Docker, Bash, Git, MongoDb, React, GraphQL)}{\href{https://eig.etalab.gouv.fr/defis/label/}{Cour de Cassation}}{Paris}{}{
  Inside the program "Entrepreneur d'Int\'{e}r\^{e}t g\'{e}n\'{e}ral", I worked in the Court of Cassation to open justice data.
  \begin{itemize}[label=\textbullet]
    \item Conducted over 40 user interviews. Weekly meeting with the Court of Cassation's executive officers to set up the roadmap.
    \item Inside a pluridisciplinary team of 4 people, I led the design and the architecture of an open source application, LABEL, to annotate and anonymize justice decisions (\emph{https://github.com/Cour-de-cassation/label})
    \item Mentoring around software quality.
    \item First MVP released in 3 months.
    \item Best software development practices: mono-repository, high test coverage, strict typing, containerization.
    \item Used by 20 public agents daily, this new tool increased the number of annotated justice decisions per user and per day from 15 to 100.
    \item Succeeded handover, the project is still under active development.
  \end{itemize}
}
\cventry{02/2019 -- 09/2020}{\href{https://360learning.com/}{Software Engineer} (JS, NodeJs, Docker, Bash, Git, MongoDb, Vue)}{\href{https://360learning.com/}{360Learning}}{Paris}{}{
  \begin{itemize}[label=\textbullet]
    \item Design, development, and production launch of the integrations to third party online courses platforms (Coursera, LinkedIn Learning, Udemy, Edx,...). Thousands of courses were synchronized daily.
    \item Redesign of the architecture of the connectors to the APIs. The time to develop a new connector for a new third party platform went down from 1 month of work to 2 days with a complete test coverage.
    \item More than 30 technical interviews realized for the the recruiting process of other software engineers.
  \end{itemize}
}
\cventry{03/2018 -- 09/2018}{\href{https://www.lip6.fr/?LANG=en}{Research Intern} (OCaml, C, Assembly, Maxima, Git)}{\href{https://www.lip6.fr/?LANG=en}{LIP6, UPMC}}{Paris}{}{
  \begin{itemize}[label=\textbullet]
    \item Automatic memory consumption analysis of program with static analysis.
    \item Implemention a functional prototype for a subset of the OCaml language (\emph{https://github.com/fondation451/ocaml\_region})
  \end{itemize}
}
\cventry{03/2017 -- 09/2017}{\href{https://ocamllabs.io/}{Research Intern} (OCaml, C, Assembly, Git)}{\href{https://ocamllabs.io/}{OCaml Labs, Cambrigde University}}{Cambrigde}{}{
  \begin{itemize}[label=\textbullet]
    \item Design and development of an open source lock-free data structure library. A lock-free data structure is a data structure which can be used safely without lock or "mutex" in a parallel environment (\emph{https://github.com/ocaml-multicore/lockfree})
  \end{itemize}
}
\cventry{06/2016 -- 08/2016}{\href{https://steep.inria.fr/en/}{Research Intern} (JS, Python, SQL, Git)}{\href{https://steep.inria.fr/en/}{STEEP, INRIA}}{Grenoble}{}{
  \begin{itemize}[label=\textbullet]
    \item Development of a web application implementing algorithms for projecting greenhouse gas consumption, constrained by different scenarios from the COP21. (\emph{http://redem.inria.fr/})
  \end{itemize}
}

\section{Education}

\cventry{2015 -- 2018}{Master's degree, Computer science (MPRI) - "parcours normalien"}{Ecole Normale Sup\'{e}rieure}{Paris}{}{
  My studies focused on functional programming, compiler design, cryptography, networking and software engineering.
}

\section{Languages}

\cvitemwithcomment{French}{Mother tongue}{}
\cvitemwithcomment{English}{Business fluent}{}
\cvitemwithcomment{Spanish}{Intermediate}{}

\end{document}
